\documentclass[12pt]{article}

\usepackage{sbc-template}
\usepackage{graphicx,url}
\usepackage[brazil]{babel}
\usepackage[utf8]{inputenc}
\usepackage[T1]{fontenc}
\usepackage{booktabs,tabularx,array,xcolor}
\usepackage[hidelinks]{hyperref}
\sloppy

% ---------------------------------------------------------------------------
% Números da base final (preenchidos a partir de base_final/log_05.json e descricao_base.json)
\newcommand{\pend}[1]{\textcolor{red}{[#1]}}
\newcommand{\NUSU}{15.133.511}
\newcommand{\NUSUtres}{2.486.475}
\newcommand{\PUSUtres}{49,0\%}
\newcommand{\DTMIN}{15/10/2010}
\newcommand{\DTMAX}{30/10/2026}
\newcommand{\TAMGB}{5,8~GB}
\newcommand{\NVAZIO}{106.169}
\newcommand{\NFUNNY}{198}
\newcommand{\NPTNULO}{94.995}
\newcommand{\NDUP}{nenhuma}
\newcommand{\DESCTXT}{Os textos são curtos: metade tem até 44 caracteres e um quarto, até 12 (amostra de 20 mil por lote). Entre os usuários, 67,5\% escreveram uma única \textit{review} AAA, e o mais ativo escreveu 1.155.}
\newcommand{\NGONULO}{635.493 (2,1\%)}
\newcommand{\LINKGIT}{\url{https://github.com/NicZorzam/projeto-final-mineracao}}
\newcommand{\LINKDRIVE}{\url{https://drive.google.com/drive/folders/1eko1pXDkzBBa40CgX2oKhJapxMtHER8U?usp=sharing}}
% ---------------------------------------------------------------------------

\title{O que os jogadores dizem sobre jogos AAA:\\ mineração de 29,6 milhões de \textit{reviews} da Steam}

\author{Nicolly Zorzam Moura\inst{1}, Pedro Jorge Câmara\inst{1}}

\address{Instituto de Computação -- Universidade Federal do Rio de Janeiro (UFRJ)\\
  Rio de Janeiro -- RJ -- Brasil
  \email{nicollyzorzammoura@gmail.com, pedrojoc@ic.ufrj.br}
}

\begin{document}
\maketitle

\begin{resumo}
\textit{Sprint} 3 -- Mineração de Dados 2026.2 (Prof. Silas Lima Filho). Este documento
contém a seção de Materiais e Métodos do artigo, incluindo a subseção de Criação da Base de Dados.
As demais seções (Introdução, Trabalhos Correlatos, Resultados etc.) são desenvolvidas nas
\textit{sprints} anteriores e seguintes.
\end{resumo}

\setcounter{section}{2}
\section{Materiais e Métodos}\label{sec:metodos}

\subsection{Propósito e questões de pesquisa}

O propósito deste trabalho é analisar \emph{o que os jogadores dizem} sobre jogos AAA nas
avaliações (\textit{reviews}) da Steam e \emph{quem são} os usuários que avaliam esses jogos.
Os trabalhos correlatos analisam \textit{reviews} da Steam com recortes por tema, comunidade ou
tipo de jogo \cite{bazzaz2026,byers2025,petrosino2022,gibrim2026}, mas nenhum isola os jogos
AAA, que concentram grande parte do volume de avaliações da plataforma. Também é raro usar o
usuário como unidade de análise; entre os trabalhos lidos, só \cite{gibrim2026} compara grupos de
usuários. Definimos três questões de pesquisa (QP):

\begin{itemize}
\item \textbf{QP1 (caracterizar):} como a recepção dos jogos AAA (taxa de recomendação, horas
jogadas até a avaliação, idioma, ano e publicadora) difere da recepção dos demais jogos da Steam?
\item \textbf{QP2 (o que dizem):} quais temas aparecem nas \textit{reviews} de jogos AAA, e
quais distinguem as avaliações negativas das positivas?
\item \textbf{QP3 (quem avalia):} que perfis de avaliadores de jogos AAA surgem ao agrupar os
usuários pelo seu histórico de avaliações, e como esses perfis se relacionam com a recomendação?
\end{itemize}

\subsection{Tipo de pesquisa}

Quanto à abordagem, a pesquisa é \textbf{quali-quantitativa}: QP1 e QP3 são respondidas com
estatística descritiva, testes de hipótese e agrupamento; QP2 combina modelagem de tópicos
(quantitativa) com a interpretação dos tópicos pelos autores (qualitativa). Quanto à natureza,
é \textbf{aplicada}, pois gera conhecimento sobre um problema concreto (a recepção de jogos AAA).
Quanto aos objetivos, é \textbf{exploratória e descritiva} \cite{gil2002}. Quanto ao
procedimento, é um \textbf{levantamento} sobre dados secundários: as \textit{reviews} já
existiam e foram analisadas sem intervenção dos pesquisadores.

\subsection{Ambiente}

A pesquisa ocorre na Steam, a loja digital de jogos para computador da Valve. Cada
\textit{review} traz uma recomendação binária (recomenda ou não recomenda), um texto livre e
metadados do autor, como as horas jogadas até a avaliação, o número de jogos na biblioteca e se o
jogo foi comprado na Steam ou recebido de graça \cite{steamapi}. A plataforma foi escolhida
porque (i) as avaliações são públicas e associadas a um identificador pseudônimo do usuário, o que
permite estudar o histórico de cada avaliador; (ii) a recomendação binária funciona como rótulo
de sentimento declarado pelo próprio autor; e (iii) é a fonte usada por todos os trabalhos
correlatos \cite{bazzaz2026,byers2025,petrosino2022,gibrim2026,jorge2023,ferreira2026,lei2025},
o que permite comparar resultados.

\subsection{Sujeitos}

Os sujeitos são os usuários da Steam que escreveram \textit{reviews} de jogos AAA. Não há
participação direta de pessoas: cada usuário aparece apenas por um identificador numérico
pseudônimo (\texttt{user\_id}), sem nome, perfil ou localização. Esse identificador é um \textit{hash} com chave (HMAC-SHA256, truncado em 63~bits) do SteamID64 do autor: é estável, o que permite agrupar as avaliações de cada usuário, mas não pode ser revertido sem a chave, que não é publicada \cite{dataset}. Como os textos são públicos e uma busca pelo texto pode levar ao perfil do autor, não reproduziremos trechos literais que permitam identificar avaliadores, seguindo os termos do \textit{dataset}. A base final reúne \NUSU{} usuários
distintos. Eles são classificados pelo número de \textit{reviews} AAA que escreveram (1, 2, 3--4,
5--9, 10--49 e 50 ou mais) e, na QP3, pelos perfis obtidos no agrupamento. Usuários com
três ou mais avaliações AAA são o foco da QP3 (\NUSUtres{} usuários, que escrevem
\PUSUtres{} das \textit{reviews} AAA).

\subsection{Criação da Base de Dados}\label{sec:base}

\subsubsection{Extração de dados}

O conjunto original é o \textit{Steam Reviews, Games and Tags}, de Sergio Del Carpio, publicado no Kaggle \cite{dataset}. Os dados foram coletados da Steam Web API e das páginas da loja por um \textit{pipeline} de código aberto de sistema de recomendação \cite{pipeline}. Usamos a versão~1, um instantâneo de 30/10/2026 (UTC); o autor publica novas versões mensalmente, com mais linhas, por isso a versão importa para a reprodução. Ela foi baixada como um arquivo compactado
(\texttt{archive.zip}, 23,3~GB compactado e 23,7~GB descompactado) com três tabelas em formato
Parquet: \texttt{games} (184.885 jogos; preço, gêneros, categorias, desenvolvedoras e
publicadoras), \texttt{game\_tags} (as 429 \textit{tags} da comunidade, com pesos, coletadas entre
29/09 e 29/10/2026) e \texttt{reviews} (121.312.346 avaliações de 2010 a 2026, em 13 arquivos de
cerca de 2~GB). As tabelas se ligam pelo \texttt{game\_id}.
Segundo a documentação, cada \textit{review} aparece uma única vez (a primeira versão coletada é mantida), as datas estão em UTC e os tempos de jogo, em minutos. Quanto aos direitos, os textos pertencem aos seus autores, e os nomes, descrições e \textit{tags} dos jogos, à Valve e às publicadoras; os dados são disponibilizados para pesquisa e educação.

A extração lê os arquivos Parquet dentro do \texttt{.zip}, sem descompactar a base inteira. Na
caracterização da \textit{sprint}~1, um leitor ``esparso'' leu o rodapé de cada Parquet e
descomprimiu apenas os blocos das colunas necessárias, o que permitiu varrer as 121,3~milhões de
avaliações numa máquina virtual com cerca de 3~GB de RAM (cerca de 50~s por arquivo). Na
\textit{sprint}~3, o mesmo leitor contou as \textit{reviews} por jogo e idioma (6~min no total), o
que permitiu comparar os critérios de seleção de jogos AAA descritos a seguir.

\textbf{Definição de AAA.} Não há definição oficial de jogo AAA, nem \textit{tag} correspondente
na Steam. Comparamos quatro critérios (Tabela~\ref{tab:criterios}): (C1) jogo publicado por uma de
17 grandes publicadoras; (C2) preço acima de um corte; (C3) a combinação dos dois; e (C4) o C1
ampliado com três publicadoras. A lista de publicadoras do C1 vem do VideoGamesCritic
\cite{vgc2026}: EA, Activision, Blizzard, Ubisoft, 2K, Rockstar, Sony, Xbox Game Studios, Bethesda,
Nintendo, Square Enix, Capcom, Bandai Namco, Sega, Warner Bros. Games, Konami e CD Projekt Red.
O site exclui ainda os jogos com a \textit{tag} Indie (C1b) e, na sua análise, usa só lançamentos a partir de 2014 com ao menos 500 avaliações; esses dois filtros não foram adotados aqui, para manter todas as avaliações dos jogos selecionados. Como o campo \texttt{game\_publishers}
tem grafias variadas (por exemplo, oito grafias da Capcom), mapeamos explicitamente 46 nomes para
16 dos 17 grupos (a Nintendo não tem jogos na Steam), mais quatro nomes para as três publicadoras
acrescentadas no C4, conferindo cada um; nomes parecidos de outras
empresas (por exemplo, ``Sony Pictures Virtual Reality'' e ``Projekt Wolfenstahl'') foram
descartados.

\begin{table}[ht]
\centering
\caption{Jogos e \textit{reviews} selecionados por critério de AAA (base completa: 120.980.379 \textit{reviews} de jogos cadastrados; 85,7\% positivas).}\label{tab:criterios}
\footnotesize
\begin{tabular}{lrrrrr}
\toprule
Critério & Jogos & \textit{Reviews} & \% base & \% positivas & PT-BR \\
\midrule
C1 -- 17 publicadoras & 1.320 & 27.319.688 & 22,6 & 84,5 & 1.522.001 \\
C1b -- 17 publicadoras, sem \textit{tag} Indie & 1.232 & 25.822.952 & 21,3 & 84,2 & 1.429.143 \\
C2a -- preço $\geq$ 39,99 & 1.872 & 16.483.479 & 13,6 & 83,9 & 675.854 \\
C2c -- preço $\geq$ 59,99 & 693 & 5.275.871 & 4,4 & 85,0 & 212.690 \\
C3 -- C1 e preço $\geq$ 39,99 & 285 & 8.984.819 & 7,4 & 84,9 & 434.764 \\
\textbf{C4 -- C1 + Koei Tecmo, Paradox e Larian} & \textbf{1.527} & \textbf{29.573.250} & \textbf{24,4} & \textbf{84,4} & \textbf{1.581.726} \\
\bottomrule
\end{tabular}
\end{table}

Adotamos o \textbf{C4}. O preço não serve como critério porque \texttt{game\_price} é o preço no
momento da coleta, e não o de lançamento: Grand Theft Auto IV aparece a 5,99, e Rainbow Six
Siege e Apex Legends, hoje gratuitos, aparecem a zero. Além disso, cortes por preço selecionam
jogos que não costumam ser chamados de AAA, como Lies of P. A moeda de \texttt{game\_price} não
consta da documentação do \textit{dataset}; os valores (por exemplo, 59,99) são compatíveis com
dólares americanos, mas isso não foi verificado.

Partimos do C1 e fizemos duas mudanças em relação à regra do VideoGamesCritic. Primeiro, não
excluímos os jogos com a \textit{tag} Indie: o recorte é pela publicadora, e a \textit{tag} é
votada pela comunidade, de modo que marca como Indie jogos publicados pelas próprias empresas da
lista, como Ori and the Will of the Wisps (Xbox Game Studios) e Risk of Rain~2 (2K). São 104 jogos
e 1,7~milhão de avaliações no C4. Segundo, acrescentamos Koei Tecmo (115 jogos, 313~mil
avaliações), Paradox (85 jogos, incluindo o selo Paradox Arc; 1,5~milhão) e Larian (8 jogos,
433~mil), porque essas três publicam jogos que, na avaliação do grupo, podem ser considerados de grande orçamento, como Nioh~2 (Koei Tecmo), Crusader Kings~III (Paradox) e Baldur's Gate~3 (Larian). Um jogo é publicado ao mesmo tempo
por Bandai Namco e Paradox.

O C4 tem limitações, que discutimos na Seção de Avaliação e Discussão: a lista ampliada combina uma
fonte externa com uma escolha do grupo; a Nintendo não tem jogos na Steam e a Blizzard tem apenas
três; e outras publicadoras de grande volume, como Valve e Krafton, ficam de fora.

\subsubsection{Tratamentos realizados}

\begin{enumerate}
\item \textbf{Filtro AAA (C4).} Mantidas apenas as \textit{reviews} dos 1.527 jogos do C4.
As 331.967 avaliações cujo \texttt{game\_id} não está cadastrado na tabela \texttt{games} ficam de
fora, porque não é possível saber a publicadora do jogo.
\item \textbf{Estouro de inteiro em \texttt{votes\_funny}.} Na base completa, 772 avaliações têm
\texttt{votes\_funny} = 4.294.967.295 ($2^{32}-1$), valor típico de um estouro de inteiro sem
sinal de 32~bits e não de votos reais. Na base final, \NFUNNY{} casos foram substituídos por nulo
e sinalizados (\texttt{votes\_funny\_overflow}), para não distorcer médias e correlações.
\item \textbf{Valores ausentes.} \NPTNULO{} avaliações AAA não têm
\texttt{playtime\_at\_review}. Mantivemos o valor nulo em vez de imputar, porque o tempo de jogo
é uma variável central da QP1 e um valor imputado criaria um padrão artificial. Criamos
\texttt{playtime\_at\_review\_h}, em horas, para facilitar a leitura. Pelo mesmo motivo, os \NGONULO{} valores ausentes de \texttt{num\_games\_owned} foram mantidos nulos; na QP3, os usuários sem esse atributo serão agrupados sem ele ou excluídos, e informaremos quantos.
\item \textbf{Textos vazios.} \NVAZIO{} avaliações têm texto vazio ou só com espaços. Elas
continuam valendo nas análises de recomendação (QP1 e QP3), mas são sinalizadas
(\texttt{texto\_vazio}) e excluídas da análise de texto (QP2). Também registramos o comprimento
do texto (\texttt{review\_len}).
\item \textbf{Duplicatas.} Verificamos se algum \texttt{review\_id} aparece mais de uma vez;
não há duplicatas.
\item \textbf{Datas de lançamento.} \texttt{game\_release\_date} é texto (por exemplo,
``Mar 27, 2018'') e inclui valores implausíveis, como 1969 e 2104, e marcadores como ``Coming
soon''. Convertemos para data, e anos anteriores a 1970 ou posteriores a 2030, assim como
marcadores, viram nulos.
\item \textbf{Contas atípicas.} Um único usuário tem 28.214 \textit{reviews} na base completa,
possivelmente um curador ou robô. Esses usuários não são removidos da base; na QP3, os acima do
percentil 99,9 de \textit{reviews} AAA são analisados à parte, para não dominarem o agrupamento.
\end{enumerate}

\subsubsection{Base final}

A base final tem \textbf{29.573.250 \textit{reviews}} de \textbf{1.527 jogos AAA} (1.481 com ao
menos uma avaliação), escritas por \NUSU{} usuários entre \DTMIN{} e \DTMAX{}, em 31 idiomas:
inglês 46,9\%, chinês simplificado 15,4\%, russo 7,7\%, português do Brasil 5,3\% (1.581.726) e
espanhol 5,2\%. A taxa de recomendação é de 84,4\%, contra 86,1\% nas 91.407.129 avaliações dos
demais jogos cadastrados. Por publicadora, ela varia de 63,4\% (Blizzard), 76,1\% (Konami) e
76,8\% (Koei Tecmo) a 92,2\% (CD Projekt Red) e 94,6\% (Larian). O número de avaliações por jogo é
muito concentrado: mediana de 3.606; o jogo com mais avaliações é Rainbow Six Siege (580.091).
\DESCTXT

A base é distribuída em Parquet (compressão zstd, \TAMGB{} no total) com duas tabelas:
\texttt{reviews\_aaa\_000} a \texttt{reviews\_aaa\_012} (uma parte por arquivo original) e
\texttt{jogos\_aaa}. A Tabela~\ref{tab:dic} resume os campos; o dicionário completo está no
repositório. A base está disponível em \LINKDRIVE{}.

\begin{table}[!ht]
\centering
\caption{Principais campos da base final.}\label{tab:dic}
\scriptsize
\begin{tabularx}{\textwidth}{>{\raggedright\arraybackslash}p{0.47\textwidth}X}
\toprule
Campo & Descrição \\
\midrule
\multicolumn{2}{l}{\textit{reviews\_aaa} (uma linha por avaliação)} \\
\texttt{review\_id, user\_id, game\_id} & identificadores da avaliação, do autor (pseudônimo) e do jogo \\
\texttt{is\_positive} & recomendação (verdadeiro = recomenda) \\
\texttt{review, review\_len, texto\_vazio} & texto livre, comprimento em caracteres e sinalizador de texto vazio \\
\texttt{language} & idioma declarado (31 valores; ex.: \textit{english}, \textit{brazilian}) \\
\texttt{reviewed\_at, updated\_at} & data da avaliação e da última edição \\
\texttt{playtime\_at\_review(\_h), playtime\_forever} & minutos (e horas) jogados até a avaliação; total jogado \\
\texttt{num\_games\_owned, num\_reviews} & jogos na biblioteca e \textit{reviews} escritas pelo autor (na Steam) \\
\texttt{votes\_up, votes\_funny, weighted\_vote\_score} & votos de útil e engraçado e pontuação de utilidade da Steam \\
\texttt{steam\_purchase, received\_for\_free} & comprado na Steam; recebido de graça \\
\texttt{written\_during\_early\_access, primarily\_steam\_deck} & escrito em acesso antecipado; jogado no Steam Deck \\
\midrule
\multicolumn{2}{l}{\textit{jogos\_aaa} (uma linha por jogo)} \\
\texttt{game\_id, game\_name, grupo\_publicadora} & jogo e grupo (uma das 17 publicadoras) \\
\texttt{release\_date, game\_price, game\_is\_free} & lançamento (tratado), preço na coleta, gratuito \\
\texttt{game\_genres, game\_categories, tag\_names, tag\_weights} & gêneros, categorias e \textit{tags} da comunidade com pesos \\
\bottomrule
\end{tabularx}
\end{table}

\subsection{Variáveis}

\textbf{QP1.} A variável dependente é a recomendação (\texttt{is\_positive}). As variáveis
independentes são: ser AAA ou não, publicadora, horas jogadas até a avaliação (em faixas: menos
de 1~h, 1--5, 5--20, 20--50, 50--100, 100--500 e mais de 500~h), idioma, ano da avaliação,
compra na Steam, jogo recebido de graça, acesso antecipado e Steam Deck.
\textbf{QP2.} O texto da avaliação é o objeto de análise, comparado entre recomendações positivas
e negativas e entre inglês e português.
\textbf{QP3.} Não há variável dependente no agrupamento. Os atributos de cada usuário são: número
de \textit{reviews} AAA, taxa de recomendação, mediana de horas até a avaliação, número de
publicadoras distintas avaliadas, fração de jogos recebidos de graça, \texttt{num\_games\_owned} e
\texttt{num\_reviews}. Depois do agrupamento, o perfil vira variável independente na análise da
recomendação.

\subsection{Procedimentos experimentais}

\textbf{QP1 -- caracterização.} Calculamos a taxa de recomendação por faixa de cada variável,
para os jogos AAA e para os demais (obtidos pela diferença entre a base completa, já agregada na
\textit{sprint}~1, e a base AAA). As diferenças de proporção são testadas com $\chi^2$ e medidas
pelo V de Cramér \cite{cramer1946}; as diferenças de horas jogadas, com o teste de Mann-Whitney
\cite{mann1947} e o delta de Cliff \cite{cliff1993}. Com dezenas de milhões de observações,
quase toda diferença é estatisticamente significativa, por isso a interpretação se apoia no
tamanho de efeito.

\textbf{QP2 -- temas.} Trabalhamos com textos em inglês (o maior grupo e o idioma da maior
parte dos trabalhos correlatos) e em português do Brasil (diálogo com os trabalhos nacionais
\cite{gibrim2026,jorge2023}). Excluímos textos vazios e com menos de 20 caracteres, que em geral não têm conteúdo
temático (por exemplo, ``good'' ou ``10/10''); pela amostra do script~06, isso deve retirar
cerca de um terço das avaliações. Depois, sorteamos uma amostra estratificada por jogo e polaridade, com até 200 avaliações por
jogo, polaridade e idioma (semente 42), para que os jogos com mais avaliações não dominem os
tópicos. Os tópicos são obtidos com o BERTopic \cite{grootendorst2022}: \textit{embeddings} de
sentenças multilíngues \cite{reimers2019,reimers2020}, redução de dimensão com UMAP
\cite{mcinnes2018} e agrupamento com HDBSCAN \cite{campello2013}. Os termos que mais distinguem
avaliações negativas de positivas são medidos com a razão de chances logarítmica com
\textit{prior} informativo \cite{monroe2008}.

\textbf{QP3 -- perfis.} Agregamos as avaliações por \texttt{user\_id} e selecionamos os usuários
com ao menos três \textit{reviews} AAA. Os atributos com cauda longa passam por $\log(1+x)$, e
todos são padronizados. O agrupamento usa o $k$-means \cite{lloyd1982}, com $k$ de 2 a 10, e os
perfis resultantes são cruzados com a recomendação. Regras de associação (Apriori
\cite{agrawal1994}) entre perfil, publicadora, faixa de horas e recomendação negativa
complementam a leitura dos grupos.

\subsection{Procedimentos de avaliação}

\textbf{QP1:} tamanho de efeito (V de Cramér e delta de Cliff), além da comparação com as taxas
de recomendação relatadas na literatura (por exemplo, 86,3\% nos jogos com geração procedural em
\cite{bazzaz2026}).
\textbf{QP2:} coerência dos tópicos ($C_V$) \cite{roder2015} e diversidade de tópicos
\cite{dieng2020}. Além disso, os dois autores leem 20 avaliações por tópico, nomeiam os tópicos de
forma independente e medem a concordância com o $\kappa$ de Cohen \cite{cohen1960}.
\textbf{QP3:} silhueta \cite{rousseeuw1987} e índice de Davies-Bouldin \cite{davies1979} para
escolher $k$; estabilidade pelo índice de Rand ajustado \cite{hubert1985} entre execuções com
subamostras de 80\%; e interpretabilidade dos perfis. As regras de associação são filtradas por
suporte, confiança e \textit{lift} $>1$.

\subsection{Material}

O processamento foi feito em Python 3.12.10, com pandas 3.1.1 \cite{mckinney2010}, PyArrow
26.0.1 e NumPy 2.5.5, num computador com Windows 10 de 64~bits (versão 10.0.19045), processador AMD Ryzen~5 2400G (3,6~GHz), 16~GB de RAM e SSD de 1~TB. A varredura da \textit{sprint}~1 rodou numa máquina virtual
Linux com cerca de 3~GB de RAM. As próximas etapas usarão scikit-learn \cite{pedregosa2011},
BERTopic, sentence-transformers, SciPy e mlxtend, com versões fixadas no arquivo
\texttt{requirements.txt} do repositório. O texto é escrito no Overleaf, com o modelo da SBC.

\textbf{Reprodutibilidade.} O repositório \LINKGIT{} contém os scripts, que devem ser executados
em ordem: \texttt{03\_reviews\_por\_jogo.py} (contagem por jogo e idioma),
\texttt{04\_criterios\_aaa.py} (Tabela~\ref{tab:criterios}), \texttt{05\_base\_final.py} (base
final e registro dos tratamentos, 6,9~min) e \texttt{06\_descrever\_base.py} (estatísticas desta seção),
além do mapeamento das publicadoras e do dicionário de dados.

\bibliographystyle{sbc}
\bibliography{referencias}

\end{document}
